\section{Typed terms and literal type support}\label{sec:syntax}

We use the lambda formulation of G\"odel's system $T$ with the argument
order of Wilken and Weiermann~\cite[Section~1]{ww}.  Its simple types are
generated by
\[
  \sigma,\tau ::= o\mid \sigma\to\tau,
  \qquad
  \lv(o)=0,\quad
  \lv(\sigma\to\tau)=\max\{\lv(\sigma)+1,\lv(\tau)\}.
\]
Arrows associate to the right.  For each type there is a countably infinite,
effectively named collection of variables of that type.  An ordinary term
is a well-typed finite expression formed from variables, abstraction,
application, and the constants
\[
\begin{aligned}
 0&:o,& S&:o\to o,\\
 D^\tau&:o\to\tau\to\tau\to\tau,&
 R^\tau&:o\to(o\to\tau\to\tau)\to\tau\to\tau.
\end{aligned}
\]
The typing rules for variables and constants use their declared types;
from $A:\sigma\to\tau$ and $B:\sigma$ one forms $AB:\tau$, and from
$A:\tau$ with $x:\sigma$ bound in its context one forms
$\lambda x:\sigma.A:\sigma\to\tau$.  Application associates to the left.
Terms are identified up to renaming bound variables.  This identification
does not include beta reduction, eta conversion, or any rearrangement of
application nodes.  Capture-avoiding substitution can be implemented by
fresh renaming or, equivalently, by a typed de Bruijn representation.

The primitive reductions are
\begin{align*}
 (\lambda x:\sigma.A)B&\longrightarrow A[x:=B],\\
 D^\tau 0\,A\,B&\longrightarrow A,
 &D^\tau(Sn)\,A\,B&\longrightarrow B,\\
 R^\tau 0\,A\,B&\longrightarrow B,
 &R^\tau(Sn)\,A\,B&\longrightarrow
      A\,n\,(R^\tau n\,A\,B).
\end{align*}
Their well-typed instances are closed under every well-typed one-hole
context, including contexts below abstractions.  In particular the first
argument of $R^\tau$ is its natural-number test, the second is its step
functional, and the third is its initial value.

\subsection{A ranked first-order sublanguage}

To state an exact recursive-path-order restriction, fix independently a
finite signature $\mathcal F$.  Each $f\in\mathcal F$ has a fixed declaration
$[o,\ldots,o]\to o$ and a fixed number of arguments.  We add the actual
ranked term former $f(A_1,\ldots,A_r):o$.  The names in $\mathcal F$ are
disjoint from the native constants and from the application symbols.
The first-order $\mathcal F$-terms are generated by these ranked formers
and variables of type $o$; they need not be closed.

For the higher-order comparison, ordinary application is represented
literally by the typed binary symbol
\[
 @_{\sigma,\tau}:[\sigma\to\tau,\sigma]\to\tau,
 \qquad
 AB=@_{\sigma,\tau}(A,B).
\]
Every native constant is a nullary head with its declared simple type.
For example $D^o0AB$ consists of a nullary $D^o$ node and three actual
application nodes.  It is not a ternary ranked case symbol.  The distinction
keeps both the native reduction rules and the first-order restriction on
their stated syntactic carriers.

All head statuses in this paper are multiset statuses.  The heads
$@_{\sigma,\tau}$ form one precedence-equivalence class, all native
constant instances form a second class strictly above it, and the heads
of $\mathcal F$ lie above both classes in a chosen strict total precedence.
This is a well-founded precedence with equivalence classes; no comparison
between native type instances is hidden in the precedence.

\subsection{The finite support used for confinement}

For a type $\sigma$, let $\operatorname{Sub}(\sigma)$ be the finite set
of its type subexpressions, including $\sigma$.  Thus
\[
 \operatorname{Sub}(o)=\{o\},\qquad
 \operatorname{Sub}(\sigma\to\tau)
 =\{\sigma\to\tau\}\cup\operatorname{Sub}(\sigma)
                         \cup\operatorname{Sub}(\tau).
\]

\begin{definition}[Literal type support]\label{def:typesupport}
The support $\Ty(A)$ of an ordinary term is the union of
$\operatorname{Sub}(\sigma)$ over the result type $\sigma$ of every
node of $A$ and the declared type of every bound variable of $A$.
Unused entries of an ambient typing context are not included.
\end{definition}

The support is finite and closed under type subexpressions.  At a ranked
ground node, the only declaration types are its actual ground argument
and result types.  In particular its absent curried type is not added to
the support.  At a native nullary head its complete original functional
type is an actual node type and is therefore included.  Auxiliary marked
and numbered terms will retain that original head declaration even when
their current output is retyped; their finite-universe condition will
check both pieces of information.

Write $\operatorname{FV}(A)$ for the finite set of free typed variable
names of $A$.  A finite type-preserving substitution $\theta$ assigns
ordinary terms of the same type to finitely many names and acts as the
identity at all other names.  Its images may contain new free names.

\begin{lemma}[Exact support under substitution]
\label{lem:typesupport-substitution}
For every ordinary term $A$ and finite type-preserving substitution $\theta$,
\[
 \Ty(A\theta)=\Ty(A)\cup
       \bigcup_{x\in\operatorname{FV}(A)}\Ty(\theta x).
\]
\end{lemma}
\begin{proof}
Induct on the literal term.  At a variable $x:\sigma$, the result type of
$\theta x$ is still $\sigma$, so
$\operatorname{Sub}(\sigma)\subseteq\Ty(\theta x)$.  This gives the
displayed equality, including the identity case.  A constant keeps its
type and has no free names.  At application and at a ranked node, the
node's result and all argument types are unchanged, and the induction
hypotheses combine by finite union.  At an abstraction, rename the binder
fresh for the finitely many relevant substitution images and mask its
name.  Its declaration and the abstraction type are retained, while
the induction hypothesis for the body contributes precisely the free
names of the abstraction.  Bound renaming changes no type support.
\end{proof}

\begin{corollary}[Support confinement is stable]
\label{cor:typesupport-stability}
Suppose $A,B$ have the same outer type,
$\Ty(B)\subseteq\Ty(A)$, and
$\operatorname{FV}(B)\subseteq\operatorname{FV}(A)$.
Then $\Ty(B\theta)\subseteq\Ty(A\theta)$ for every such $\theta$.
Moreover $\Ty(C[B])\subseteq\Ty(C[A])$ for every well-typed one-hole
context $C$ of the appropriate hole type.
\end{corollary}
\begin{proof}
The substitution assertion follows by inclusion in both unions of
Lemma~\ref{lem:typesupport-substitution}.  A fixed context contributes
the same finite set of node and binder types to both plugged terms;
the two hole types agree.  Taking the union with those common types
preserves support inclusion.
\end{proof}

This support condition permits one finite type universe to control every
descendant of a fixed source.  It does not bound the size of descendants,
the length of generated argument lists, or the types that may occur in
other, unrelated input terms.
